\documentclass[a4paper,12pt]{article}

\usepackage[T2A]{fontenc}
\usepackage[utf8]{inputenc}
\usepackage[russian]{babel}

\usepackage{amsmath, amssymb}
\usepackage{geometry}
\usepackage{fancyhdr}
\usepackage{enumitem}

\geometry{margin=2cm}
\setlength{\parskip}{6pt}
\setlength{\parindent}{0pt}

\pagestyle{fancy}
\fancyhf{}
\rhead{7 класс}
\cfoot{\thepage}

\begin{document}

\begin{center}
    {\Large \textbf{Делимость.}}\\[0.5em]
\end{center}
    
\begin{enumerate}[label=\textbf{Задача \arabic*.}]
\item Площадь прямоугольника с натуральными сторонами равна 57. Чему равен периметр фигуры?
\item Чему равно $m - n$ для натуральных n, m, если известно, что $m^2 - n^2$ - простое число?
\item Найдите наименьшее натуральное число n, такое что
n! делится на 990 ($n! = 1 \cdot2\cdot3\cdot\ldots\cdot n$).
\item Известно, что дробь $\frac{a}{b}$ - сократима. Верно ли, что сократима дробь $\frac{a-b}{a+b}$?
\item Сумма двух натуральных чисел равна 177. Докажите, что произведение этих чисел не может делиться на 177.
\item Найдите двузначное число, первая цифра которого равна разности между этим числом и числом, записанным теми же цифрами, но в обратном порядке.
\item В ряд выписаны числа от 1 до 25. Можно ли расставить между ними знаки «+» и «–» так, чтобы
значение полученного выражения было равно нулю?
\item При каких натуральных n выражение $n^3 - n$ кратно 6?
\item Докажите, что число $3^{100} - 1$ делится на 2.
\item Числа $1, 2, 3, …, 456$ выписаны в строчку по порядку. За ход разрешается менять местами любые два числа, стоящие через одно (например, 196 и 198). Можно ли с помощью таких ходов расположить все числа в обратном порядке? 
\item Найдите остаток числа $26^{36}$ при делении на 7.
\item Какая натуральная степень числа 2 оканчивается четырьмя одинаковыми цифрами?
\item Докажите, что для всякого натурального числа n $7^n +5$ кратно 6.
\item Воспользуясь определением факториала ($n! = 1 \cdot2\cdot3\cdot\ldots\cdot n$), покажите, что существует такое натуральное число n, что все числа $n + 1, n + 2, . . . , n + 2025$ -  составные.
\item Докажите, что число является квадратом целого числа тогда и только тогда, когда оно имеет нечетное число натуральных делителей.
\item Верно ли, что бесконечных чисел бесконечно много?
\item Существуют ли 100 подряд идущих составных чисел?
\item Докажите, что для любого n существует n подряд идущих составных чисел (то есть существует сколь угодно длинная цепочка составных чисел, соседи в которой отличаются на 1).
\end{enumerate}
\end{document}